\documentclass[11pt]{article}
\usepackage[T1]{fontenc}
\usepackage{lmodern,microtype,amsmath,amssymb,amsthm,mathtools}
\usepackage[margin=29mm]{geometry}
\usepackage{booktabs,array,enumitem}
\usepackage{xcolor}
\usepackage[colorlinks=true,linkcolor=blue!55!black,citecolor=blue!55!black,urlcolor=blue!55!black]{hyperref}
\setlength{\emergencystretch}{2em}
\numberwithin{equation}{section}
\newtheorem{theorem}{Theorem}[section]
\newtheorem{proposition}[theorem]{Proposition}
\newtheorem{lemma}[theorem]{Lemma}
\newtheorem{corollary}[theorem]{Corollary}
\newtheorem{conjecture}[theorem]{Conjecture}
\theoremstyle{definition}
\newtheorem{definition}[theorem]{Definition}
\theoremstyle{remark}
\newtheorem{remark}[theorem]{Remark}
\newcommand{\Hhyp}{H_{\mathrm{hyp}}}
\newcommand{\R}{\mathbb R}
\newcommand{\dd}{\,\mathrm d}
\DeclareMathOperator{\degform}{deg}
\hypersetup{pdftitle={Mechanical renewal and D²(ln D)² lower bounds for Hilbert numbers},
 pdfauthor={Haibo Lu},pdfsubject={Mechanical construction of hyperbolic limit cycles},
 pdfkeywords={Hilbert number, polynomial vector field, Abelian integral, mechanical renewal}}
\title{Mechanical renewal and $D^2(\ln D)^2$\\ lower bounds for Hilbert numbers}
\author{Haibo Lu\\\small Shanghai Institute of Technology, Shanghai 201418, China\\
  \small\texttt{luhaibo1985@gmail.com}}
\date{24 October 2026}
\begin{document}
\maketitle
\begin{abstract}
For every sufficiently large integer $D$, we construct a real planar
polynomial vector field of degree at most $D$ with at least
$D^2(\ln D)^2/[8(\ln 2)^2]-O(D^2\ln D)$ hyperbolic limit cycles.
Qin and Sun previously obtained the order $D^2(\ln D)^2$ through two
3-adic filtrations on a Chebyshev center grid. We give a different construction based on
auxiliary mechanical folds and planar fourfold renewal, improving the leading
coefficient even after polynomial replication extends their subsequence
coefficient to all sufficiently large degrees.
The auxiliary folds contribute a logarithmic number of birth levels, while
planar renewal supplies the second logarithm. Reflections cancel the inherited and opposite-axis periods on
each new family, allowing all selected zeros to be realized by one polynomial
perturbation. A triangular allocation of monomials closes the total-degree
budget, while small even terms fill the intervening degrees. We give an
explicit finite lower bound.
\end{abstract}

\section{The problem and the construction}

For an integer $D\ge1$, consider real planar polynomial vector fields
\begin{equation}\label{eq:field}
 \dot x=P(x,y),\qquad \dot y=Q(x,y),\qquad
 \max(\deg P,\deg Q)\le D.
\end{equation}
Let $H(D)$ be the supremum of the cardinalities of finite families of
distinct limit cycles carried by a single field \eqref{eq:field}, and
let $\Hhyp(D)$ be the same supremum counting only hyperbolic cycles.
Thus $H(D)\ge\Hhyp(D)$. We ask how many cycles can be obtained at a fixed total-degree
budget when an old configuration is copied and new cycles are inserted
at each stage.

Christopher--Lloyd \cite[Section~3]{ChristopherLloyd1995} introduced
fourfold quadratic replication followed by new cycle nests near axis
centers to obtain bounds of order $D^2\ln D$. Li--Chan--Chung
\cite[Section~3, Remark~4 and Proposition~5]{LiChanChung2002}
corrected the axial perturbations in that argument. Itenberg--Shustin
\cite[Theorem~2.1]{ItenbergShustin2000} obtained an all-degree
$D^2\ln D$ lower bound by a Viro-type gluing construction. Han--Li
\cite[Lemma~2.6 and Theorem~2.1]{HanLi2012} developed a two-axis
recursion: perturbing the second component by $x$ times an odd
polynomial in $y$ creates nests near the nonzero $x$-axis centers;
a smaller odd perturbation in the first component then creates nests
near the $y$-axis centers while retaining the first nests. Their
Theorem~1.3 gives the all-degree result
\[
 \liminf_{D\to\infty}\frac{H(D)}{D^2\ln D}\ge\frac1{2\ln2}.
\]
Their recurrence is stated for odd-multiplicity cycles, but its births
in Lemma~2.6 are hyperbolic. Starting with no selected cycles in the
Hamiltonian seeds of their Lemma~3.3, each finite iteration preserves
hyperbolicity under the fourfold pullback and the subsequent small
perturbations. Thus the same liminf bound holds for $\Hhyp$.

Related mechanical ingredients also have direct precedents:
Li--Li--Llibre--Zhang
\cite[p.~445]{LiLiLlibreZhang2001} count two-index monomial periods
under a triangular degree constraint for a homogeneous potential;
Han--Romanovski \cite[Theorems~2.1 and 3.2]{HanRomanovski2011} combine
one-coordinate doubling with new outer cycles in polynomial
Li\'enard systems, already producing one logarithmic accumulation.

The two-source renewal in \cite{Lu2026} gives
$\Hhyp(D)\ge D^2\log_2D-O(D^2)$, where the subscript denotes the base.
Qin--Sun \cite[Theorem~1.1 and Section~6.2]{QinSun2026} subsequently obtained a lower
bound of order $D^2(\ln D)^2$. For every
$r\ge2$ they construct a field of degree at most
\begin{equation}\label{eq:QS}
 d_r=4\cdot3^r-5
 \quad\text{with at least}\quad
 \ell_r=4r^2 3^{2r-2}-(3^r-1)^2
\end{equation}
hyperbolic limit cycles. The coefficient of $d_r^2(\ln d_r)^2$ is
$c_{\mathrm{QS}}=1/[36(\ln3)^2]$. Their hierarchical Melnikov construction combines two
3-adic filtrations, a local Cartesian-rank argument under sufficiently
strong anisotropy, and a simultaneous choice of perturbation coefficients.
Within the fourfold copying framework, we construct separable Hamiltonians
and perturbing forms directly. We allocate the mechanical monomials
according to their degrees after all remaining folds and use the resulting
auxiliary history as a source at every planar renewal.
This gives the same logarithmic order by a different source mechanism
and improves the coefficient. The comparison with the two-source renewal
also involves its input class: that step accepts an arbitrary selected
hyperbolic configuration, whereas the present construction uses a
specially built sequence of separable Hamiltonians. Section~\ref{sec:replication}
compares both the mechanisms and the resulting bounds.

Our main result holds for every sufficiently large integer degree.
Throughout the paper, $\ln$ denotes the natural logarithm.

\begin{theorem}\label{thm:main}
There are absolute constants $C>0$ and $D_0$ such that, for every
integer $D\ge D_0$, there exists a real planar polynomial vector field
of degree at most $D$ having at least
\[
 \frac{D^2(\ln D)^2}{8(\ln2)^2}-CD^2\ln D
\]
distinct hyperbolic limit cycles. In particular,
\begin{equation}\label{eq:asymptotic}
 \Hhyp(D)\ge
 \frac{D^2(\ln D)^2}{8(\ln2)^2}-O(D^2\ln D).
\end{equation}
\end{theorem}

Proposition~\ref{prop:mechanical-certificate} gives an explicit finite lower bound
$L(n)$ for every integer $n\ge2$, with field degree at most $2n-1$.
On the dyadic sequence it simplifies to
\begin{equation}\label{eq:dyadic}
 L(2^r)=\frac{4^r r(r+1)}2-\frac83\,4^r+4\cdot2^r-\frac43,
 \qquad r\ge1.
\end{equation}
For $r=2,3,4,5$, this gives at least $20,244,1940,12756$
cycles at degrees $7,15,31,63$.

The quantitative increment over \cite[Proposition~4.1]{Lu2026} comes
from the source inserted at each stage. That renewal has degree bound $m^+\le2m+5$
and count $N^+=4N+4(m+2)(m+1)$, whose normalized birth term stays of
constant order. Here a mechanical Hamiltonian $z^2/2+V(w)$ undergoes
its own sequence of $w$ folds: each fold doubles the earlier period
zeros and creates a new outer family. We allow both the kinetic-coordinate exponent
and the potential-variable exponent to vary. After accounting for the
remaining real copies, each auxiliary birth level contributes at
quadratic order within the final source budget.

In the planar Hamiltonian $U(x)-V(y)$, the two-coordinate fold then
produces four old copies and inserts this source near both axes.
On the dyadic sequence, let $n_i=2^i$ and let $K_i$ denote the number
of newly selected period zeros at stage $i$. Its $i$ auxiliary birth levels give
$K_i=i n_i^2-4(2^i-1)$. If $L_i$ is the cumulative count, then $L_i=4L_{i-1}+K_i$, and
the normalized count $L_i/n_i^2$ increases by
$i+O(2^{-i})$, rather than by a constant, and summing over $i$ produces
the second logarithm. The extra history supplies this factor $i$;
the fold-weighted triangular allocation places the source within the
Hamiltonian field degree $2n_i-1$ and yields the coefficient in
\eqref{eq:asymptotic}.

There are two obstacles to this counting argument. First, the new sources
must coexist with the old ones in a single polynomial form. On every
new axis oval a reflection reverses orientation while preserving both
the old pullback form and the opposite-axis source. Their periods vanish
identically there. The two new amplitudes can therefore be chosen small
enough to preserve all old simple zeros without altering the desired new
zeros. Second, the two uses of folding must fit the same total-degree
budget. A mode born $k$ auxiliary folds earlier has degree in the potential variable
$2^k(2q+1)-1$; allocating the kinetic-coordinate exponent against this exact cost
closes the budget. To cover every integer degree we add small even
terms to the potentials. We fix all future auxiliary budgets at the
outset, so only finitely many already chosen zeros must survive each
addition.

For every integer $D\ge31$, Qin--Sun's direct all-degree corollary gives
a lower bound with leading coefficient $c_{\mathrm{QS}}/9$ and retains
the seed hyperbolicity
\cite[Corollary~1.2 and Sections~6.2, 6.4]{QinSun2026}. Section~\ref{sec:replication}
uses the replication of Eshkobilov--Kadyrov--Mamayusupov \cite{EKM2026}
to recover $c_{\mathrm{QS}}$ for all sufficiently large degrees, with
remainder $O(D^2\ln D\ln\ln D)$. Thus the ratio of the leading
coefficients of the present construction and this all-degree consequence is
\[
 \frac{1/[8(\ln2)^2]}{1/[36(\ln3)^2]}
 =\frac92\left(\frac{\ln3}{\ln2}\right)^2\approx11.3045.
\]

\section{Simple period zeros and finite simultaneous choices}\label{sec:prelim}

We use polynomial one-forms $\omega=A\,dx+B\,dy$, with
$\deg\omega=\max(\deg A,\deg B)$. A period zero always refers to a
specified connected family of regular compact ovals, parametrized by the
Hamiltonian value. Distinct ovals can have the same value. We select
pairwise disjoint small closed annuli about the finitely many ovals that
will produce cycles.

\begin{lemma}\label{lem:jets}
Let $f_1,\ldots,f_M$ be linearly independent real analytic functions on a
connected open interval $I$. For every nonempty open subinterval $J\subset I$
there is a real linear combination with at least $M-1$ distinct simple
zeros in $J$.
\end{lemma}
\begin{proof}
For analytic functions linear independence implies that the Wronskian is
not identically zero. One way to see this is to perform elimination on
their convergent Taylor series at an interior point: independent analytic
germs admit a basis with strictly increasing orders of vanishing, whose
Wronskian has a nonzero leading term. The identity theorem transfers
independence to every open subinterval. Choose a point $h_0$ there where
the Wronskian is nonzero. The map from coefficients to the Taylor
coefficients through order $M-1$ at $h_0$ is invertible.

Fix distinct real numbers $t_1,\ldots,t_{M-1}$ and prescribe those Taylor
coefficients to be the coefficients of
$\prod_{j=1}^{M-1}(h-h_0-\varepsilon t_j)$. The resulting coefficient
vector stays bounded as $\varepsilon\to0$. For the corresponding
function $F_\varepsilon$,
\[
 \varepsilon^{-(M-1)}F_\varepsilon(h_0+\varepsilon t)
 =\prod_{j=1}^{M-1}(t-t_j)+O(\varepsilon)
\]
in $C^1$ on each fixed compact $t$-interval. Each simple root persists,
and its energy lies in the prescribed subinterval for sufficiently small
$\varepsilon>0$. For $M=1$ the assertion is empty.
\end{proof}

\begin{lemma}\label{lem:persistence}
For a finite family of regular compact Hamiltonian annuli and fixed
smooth one-forms, selected simple period zeros persist under a sufficiently
small smooth change of the Hamiltonian and forms. A single smallness
threshold can be used for the whole finite family.
\end{lemma}
\begin{proof}
On each compact annulus choose a transverse coordinate and smoothly
parametrize nearby level curves; regularity gives this parametrization
by the implicit function theorem. Differentiating the integral in the
energy parameter shows $C^1$ convergence of the period functions when
the Hamiltonian and forms converge in sufficiently high finite
differentiability norms on a neighborhood of the annulus.
For each simple zero choose a closed energy interval $J$ on which
$|M'|\ge a_J>0$ and the two endpoint values have opposite signs and
absolute value at least $b_J>0$. Taking the minimum of these finitely
many positive margins preserves both properties under a common
$C^1$ perturbation. The intermediate value theorem and monotonicity
give a simple zero in each selected interval. The annuli remain disjoint
under a sufficiently small perturbation.
\end{proof}

The common threshold in this lemma is taken only after the degree and
all the finite forms under consideration have been fixed. No threshold
uniform in the degree is used.

The following first-order criterion is standard; see, for example,
\cite[p.~443]{LiLiLlibreZhang2001}. We record its finite simultaneous form.

\begin{lemma}\label{lem:melnikov}
Suppose that a polynomial Hamiltonian $H$ and a polynomial one-form
$\omega=A\,dx+B\,dy$ have $N$ selected simple period zeros in pairwise
disjoint compact regular annuli. For all sufficiently small nonzero real
$\mu$, the single polynomial field
\begin{equation}\label{eq:singlefield}
 X_\mu=(H_y-\mu B,\,-H_x+\mu A)
\end{equation}
has at least $N$ distinct hyperbolic limit cycles and degree at most
$\max(\deg H-1,\deg\omega)$.
\end{lemma}
\begin{proof}
In the energy coordinate on a local transversal the return displacement
is $\mu M(h)+O(\mu^2)$ in $C^1$, where
$M(h)=\oint_{H=h}\omega$ is oriented by the Hamiltonian flow.
This expansion follows by smooth dependence of the flow and return time
and differentiation at $\mu=0$. Apply Lemma~\ref{lem:persistence} to the
displacement divided by $\mu$, using the minimum of the finite margins.
At each resulting fixed point the derivative of the return map is
$1+\mu M'(h)+O(\mu^2)$ and differs from $1$ for sufficiently small
nonzero $\mu$. It is positive and close to $1$, so the cycle is
hyperbolic. The degree follows directly from \eqref{eq:singlefield}.
\end{proof}

\section{A mechanical construction with a finite auxiliary history}
\label{sec:mechanical}

The construction uses an indefinite separable Hamiltonian.  Two new
families of compact ovals occur near the coordinate axes after each
quadratic fold.  The periods of each source form vanish on the opposite family by
reflection.  The additional logarithm comes from an auxiliary mechanical
history: at a given global stage, all earlier auxiliary birth levels
contribute at the same leading order after their real copies are counted.
For arbitrary degrees, the potentials are padded by small positive even
monomials.  We fix the complete finite list of future auxiliary budgets
before making these perturbations.

The period calculation uses the monomial powers of
\cite[p.~445]{LiLiLlibreZhang2001} and the outer-period asymptotics of
\cite[Lemma~2.3]{HanRomanovski2011}. We give the required version for
arbitrary even potentials, since each fold changes the potential.

\begin{lemma}[Mechanical outer periods]\label{lem:mechanical-rank}
Let $W$ be an even real polynomial of degree $2s$ with positive leading
coefficient.  On the outer period annulus of
$F(z,w)=z^2/2+W(w)$, the periods
\[
 I_{p,q}(h)=\oint_{F=h}z^{2p+1}w^{2q}\,dw,
 \qquad p\geq0,\quad 0\leq q<s,
\]
are linearly independent for every finite collection of distinct pairs
$(p,q)$.  A collection of $M$ such forms therefore has a linear
combination with $M-1$ simple period zeros in any prescribed nonempty
open interval of the outer annulus.
\end{lemma}

\begin{proof}
Write the leading term of $W$ as $c w^{2s}$, $c>0$.  Scaling $w$ by
$(h/c)^{1/(2s)}$ in the integral between the two outer turning points gives
\[
 I_{p,q}(h)\sim C_{p,q}
 h^{p+1/2+(2q+1)/(2s)},\qquad C_{p,q}\ne0.
\]
For a consistent oval orientation the constants have a fixed sign:
apart from that sign, the scaled integral is a positive multiple of
$\int_{-1}^{1}(1-u^{2s})^{p+1/2}u^{2q}\,du$.
The exponents are distinct, since their equality would imply
$p-p'=-(q-q')/s$ with $|q-q'|<s$.  The term with greatest exponent
excludes every nontrivial finite linear relation.  The periods are
analytic on the connected outer energy interval, so their independence
holds on each open subinterval.  Lemma~\ref{lem:jets} gives the stated
simple zeros.  Each zero can be assigned a compact regular annular
neighborhood and a positive $C^1$ persistence margin.
\end{proof}

Reflection cancellation of inherited periods is already used in the
classical fourfold construction \cite[equation~(3.3)]{ChristopherLloyd1995}.
The following observation records its persistence under the small
Hamiltonian changes needed here.

\begin{lemma}[Symmetry on continued ovals]\label{lem:continued-symmetry}
Let $F$ be a smooth planar Hamiltonian satisfying $F\circ R=F$ on a
neighborhood of a compact regular annular family, where $R$ is a
coordinate reflection preserving each oval of the family.
If $G\circ R=G$ and $G$ is sufficiently close to $F$ in $C^3$ on a
slightly larger annulus, the nearby regular ovals of $G$ are also
preserved by $R$.  If $R^*\alpha=\alpha$, the period of $\alpha$ on
this entire nearby family is identically zero.
\end{lemma}

\begin{proof}
Use the flow of $\nabla F/|\nabla F|^2$ to write an invariant annular
collar as $(\theta,h)$ with $F=h$.  The flow is equivariant under $R$.
For $G$ close to $F$, the equation $G(\theta,u)=h$ has a unique solution
$u=u(\theta,h)$ in this collar, by the implicit function theorem.
Uniqueness and $G\circ R=G$ make its closed normal graph invariant under
$R$.  It is a regular closed level component, hence an actual periodic
oval.  A reflection reverses its oval orientation, and consequently
\[
 \oint\alpha=\oint R^*\alpha=-\oint\alpha=0.
\]
This holds for every energy in an interval, so it includes the energy
derivative of the period.  The same collar construction gives the
$C^1$ continuity of periods used below, as in
Lemma~\ref{lem:persistence}.
\end{proof}

\begin{proposition}[Mechanical renewal]\label{prop:mechanical-certificate}
Fix an integer $n\geq2$ and put
\[
 \begin{gathered}
 r=\lfloor\log_2 n\rfloor,\qquad
 n_i=\left\lfloor\frac{n}{2^{r-i}}\right\rfloor
 \quad(0\leq i\leq r),\\
 \eta_i=n_i-2n_{i-1}\in\{0,1\}\quad(1\le i\le r).
 \end{gathered}
\]
For $1\leq i\leq r$ and $0\leq k<i$, set
$e_{i,k}=n_{i-1}\bmod 2^k$ and
\begin{equation}\label{eq:mechanical-birth}
 K_i=i(n_i^2-\eta_i)
       -4\sum_{k=0}^{i-1}(e_{i,k}^2+\eta_i e_{i,k})
       -4(2^i-1).
\end{equation}
There are even polynomials $U_r,V_r$ of degree $2n$ with positive
leading coefficients and a polynomial one-form $\omega_r$ of degree
at most $2n-1$ such that its period functions on compact regular annuli
of $H_r(x,y)=U_r(x)-V_r(y)$ have at least
\begin{equation}\label{eq:mechanical-total}
 L(n)=\sum_{i=1}^{r}4^{r-i}K_i
\end{equation}
simple zeros with pairwise disjoint annular neighborhoods.
In particular, there exists a real polynomial vector field of degree at
most $2n-1$ having at least $L(n)$ hyperbolic limit cycles. The bound satisfies
\begin{equation}\label{eq:finitebound}
 L(n)\ge\frac12n^2r^2+\frac16n^2r-\frac{22}{3}n^2.
\end{equation}
\end{proposition}

\begin{proof}
We construct the Hamiltonian and form by a finite induction, retaining
the auxiliary forms needed at later stages. The estimate
\eqref{eq:finitebound} is proved in Section~\ref{sec:counting}.

\paragraph{Auxiliary budgets.}
For each future global stage $t\in\{1,\ldots,r\}$, we build an auxiliary
form through levels $j=0,\ldots,t-1$ for each of the potentials $U_j$
and $V_j$.  Denote either such form by
\[
 \Omega_j^{(t)}=\sum_{p\geq0}z^{2p+1}R_p(w)\,dw.
\]
The final pullback into a stage-$t$ axis source has degree at most
$2(p+\deg R_p+1)$, so its endpoint budget is
$p+\deg R_p\leq n_t-2$.  Reversing the remaining folds gives the
budget at level $j$:
\begin{equation}\label{eq:mechanical-future-budget}
 p+2^{t-1-j}(\deg R_p+1)-1\leq n_t-2
 \quad\text{for each nonzero }R_p.
\end{equation}
A fold $w\mapsto w^2-b$ replaces $R_p$ by
$2wR_p(w^2-b)$, with
$\deg(2wR_p(w^2-b))+1=2(\deg R_p+1)$ for $R_p\ne0$.
Thus it preserves~\eqref{eq:mechanical-future-budget} at
the next level.  A birth at level $\ell$ uses exactly the monomials
\begin{equation}\label{eq:mechanical-birth-modes}
 z^{2p+1}w^{2q}\,dw,\qquad
 0\leq q<n_\ell,\qquad
 0\leq p\leq n_t-1-2^{t-1-\ell}(2q+1).
\end{equation}
They satisfy the budget.  Their periods are independent by
Lemma~\ref{lem:mechanical-rank}.

Set $U_0(w)=V_0(w)=w^2$ and $\omega_0=0$.  For every label $t$,
choose the level-zero auxiliary form in the span of the monomials
\eqref{eq:mechanical-birth-modes}, with one fewer selected simple period
zero than the dimension of that span, and with compact regular annular
neighborhoods around those zeros.  For a one-dimensional
space this requirement is empty.  The complete list of labels is fixed
at this point.

Suppose that stage $i-1$ has been completed.  Its data are the actual
potentials, one global form and all its selected zeros, and the
auxiliary forms $\Omega_{i-1}^{(t)}$ for $t\geq i$, separately for $U$
and $V$.  Every selected zero has a compact regular annular neighborhood
and a positive $C^1$ margin.  The list is finite.

\paragraph{Real copies and the two mechanical sources.}
Write $U=U_{i-1}$, $V=V_{i-1}$, and first consider
\[
 \widehat H_i(x,y)=U(x^2-a)-V(y^2-b),\qquad
 \Phi(x,y)=(x^2-a,y^2-b).
\]
For sufficiently large positive $a,b$, the four real inverse branches
of $\Phi$ cover all selected old global annuli.  They give four
diffeomorphic copies of each old simple period zero for
$\Phi^*\omega_{i-1}$, up to orientation.  The maps $(z,w)\mapsto(z,w^2-a)$ and
$(z,w)\mapsto(z,w^2-b)$ likewise have two diffeomorphic inverse branches
on the selected $U$- and $V$-auxiliary annuli, respectively. Pulling back
the auxiliary forms preserves their period functions up to orientation,
so every selected auxiliary zero has two simple copies.

The current vertical source uses the already fixed $V$-auxiliary form
$\Omega_{i-1,V}^{(i)}$.  Put
\[
 c_a=\sqrt{2U'(a)},\qquad z=c_a x,\qquad w=y^2-b.
\]
On either real $y$ branch the Hamiltonian deficit is exactly
\begin{equation}\label{eq:mechanical-limit}
 U(a)-\widehat H_i=K_a(z)+V(w),\qquad
 K_a(z)=U(a)-U\left(-a+\frac{z^2}{2U'(a)}\right).
\end{equation}
Since $U$ is even with positive leading coefficient, $U'(a)>0$ for
large $a$.  Taylor expansion gives $K_a(z)=z^2/2+o(1)$ in every fixed
$C^j$ norm on compact $z$ intervals.  More precisely, if $d=\deg U$,
the coefficient of $z^{2k}$ for $k\geq2$ is
$O(a^{-(k-1)d})$. Every regular mechanical oval is individually fixed
by $z\mapsto-z$: over its $w$ interval it consists of the two graphs
$z=\pm\sqrt{2(h-V(w))}$ joined at their turning points. This holds
also for a potential well centered at a nonzero value of $w$. The
$C^j$ convergence of $K_a$ to $z^2/2$ gives continuations of the auxiliary
ovals and their simple period zeros to \eqref{eq:mechanical-limit}, by
Lemma~\ref{lem:persistence}. Lemma~\ref{lem:continued-symmetry} preserves
the reflection symmetry of these continued ovals.  Pull back the fixed auxiliary form to
the polynomial source
\begin{equation}\label{eq:mechanical-vertical}
 \alpha_V=
 \sum_p(c_a x)^{2p+1}R_p(y^2-b)\,2y\,dy.
\end{equation}
Its periods on the two branches are $C^1$ close, up to orientation, to
the selected auxiliary periods.  This produces two copies of every
auxiliary simple zero.

Interchanging $U,x,a$ with $V,y,b$ gives a horizontal source
\begin{equation}\label{eq:mechanical-horizontal}
 \alpha_U=
 \sum_p(c_b y)^{2p+1}\widetilde R_p(x^2-a)\,2x\,dx,
 \qquad c_b=\sqrt{2V'(b)},
\end{equation}
with two copies of the corresponding $U$-auxiliary zeros.
The vertical approximation uses only large $a$ and fixed previous
$V$-annuli; $b$ only has to cover their $w$ ranges.  The horizontal
approximation uses large $b$ and fixed previous $U$-annuli, while $a$
covers their $w$ ranges.  Thus these are compatible lower bounds on
$a,b$.  In particular, no approximation is required near the new outer
barriers created by the present folds.

All these annuli can be made disjoint.  Old global copies have
$|x|\sim\sqrt a$, $|y|\sim\sqrt b$; the vertical sources have
$x=O(1/c_a)$ and $|y|\sim\sqrt b$; the horizontal sources have
$y=O(1/c_b)$ and $|x|\sim\sqrt a$.  Within a source, disjointness
comes from the selected auxiliary annular neighborhoods.  The ovals
are genuine closed level components, obtained in regular annular
collars; properness of the indefinite Hamiltonian is not needed.

\paragraph{Positive even padding.}
Fix such $a,b$, and define
\begin{equation}\label{eq:mechanical-padded-potentials}
 U_i(x)=U(x^2-a)+\epsilon x^{2n_i},\qquad
 V_i(y)=V(y^2-b)+\delta y^{2n_i},\qquad H_i=U_i-V_i.
\end{equation}
If $\eta_i=0$, set $\epsilon=\delta=0$.  If $\eta_i=1$, choose both
coefficients strictly positive and sufficiently small.  The resulting
potentials have exact degree $2n_i$, are even, and have positive leading
coefficients.

The finite conditions on this choice are the persistence of the four
old global copies, the two auxiliary copies for every label $t>i$,
and the current vertical and horizontal source zeros for the fixed
forms~\eqref{eq:mechanical-vertical} and
\eqref{eq:mechanical-horizontal}.  All are conditions on already fixed
compact sets and forms.  In vertical source coordinates the new error
is exactly
\[
 -\epsilon(z/c_a)^{2n_i}+\delta(w+b)^{n_i}.
\]
For fixed $a,b,c_a$ its $C^3$ norm on each selected compact rectangle
tends to zero with $(\epsilon,\delta)$.  In horizontal coordinates the
error is $\epsilon(w+a)^{n_i}-\delta(z/c_b)^{2n_i}$ and has the same
property.  The auxiliary and old global errors also tend to zero on
their fixed collars.  Lemma~\ref{lem:persistence} therefore imposes
only finitely many compatible upper bounds on the positive pair
$(\epsilon,\delta)$.

The continued vertical ovals are invariant under
$R_x(x,y)=(-x,y)$, and the horizontal ovals under
$R_y(x,y)=(x,-y)$, by Lemma~\ref{lem:continued-symmetry}.  Both
reflections leave $H_i$ invariant.  Functoriality of pullback gives
\[
 R_x^*\Phi^*\omega_{i-1}=R_y^*\Phi^*\omega_{i-1}
 =\Phi^*\omega_{i-1},\qquad
 R_x^*\alpha_U=\alpha_U,\quad R_y^*\alpha_V=\alpha_V.
\]
Consequently, on the entire vertical annuli the periods of
$\Phi^*\omega_{i-1}$ and $\alpha_U$ vanish exactly.  On the entire
horizontal annuli the periods of $\Phi^*\omega_{i-1}$ and $\alpha_V$
vanish exactly.  Padding changes the old copied periods by a small
amount, but does not change any of these fresh-source identities.

\paragraph{Auxiliary births and the global form.}
With the padded potentials now fixed, complete the auxiliary forms for
each remaining label $t>i$.  Start with the one-coordinate pullback of
$\Omega_{i-1}^{(t)}$, whose selected simple zeros have survived.  On a
new outer annulus of $z^2/2+U_i(w)$ or $z^2/2+V_i(w)$, this pulled form
has identically zero period: it is invariant under $w\mapsto-w$,
which reverses the oval orientation.  In the span of \eqref{eq:mechanical-birth-modes}, choose a fresh birth
form with one fewer selected simple period zero than the dimension of
that span. Add it to the pulled form with one sufficiently small nonzero
scalar. The sum retains every selected inherited zero and supplies the
chosen new simple zeros on an outer interval.
For either padded potential $W=U_i$ or $V_i$, choose the new outer
interval above all real critical values of $W$ and all values of
$z^2/2+W(w)$ on the previously selected compact annular neighborhoods.
Such an interval exists: these neighborhoods are finite in number and
compact, and $W$ has positive leading coefficient. Each mechanical level
there is a unique outer oval, so the new annular neighborhoods are disjoint
from those previously selected. The degree $2n_i$ and positive leading
coefficient give the required rank, however small the padding coefficient.
The forms for different future stages are chosen independently; these
choices change neither $H_i$ nor the current global form.

Finally define the single global form
\begin{equation}\label{eq:mechanical-global-form}
 \omega_i=\Phi^*\omega_{i-1}
                  +\lambda_V\alpha_V+\lambda_U\alpha_U.
\end{equation}
Choose two nonzero scalars sufficiently small to retain all four old
global copies.  On a fresh vertical annulus its period is exactly
$\lambda_V\oint\alpha_V$, and on a fresh horizontal annulus it is
exactly $\lambda_U\oint\alpha_U$.  Thus all new zeros coexist for this
same coefficient vector.  Each newly chosen scalar is fixed before
the next stage: the resulting positive, possibly very small, simple-zero
margins are added to the induction data.  No uniform margin over the
future degree or over all polynomial budgets is required.  This completes the finite induction.

\paragraph{Degrees and number of zeros.}
For an auxiliary monomial with $w$-coefficient degree $d$, the vertical
or horizontal source degree is at most
\[
 (2p+1)+2d+1=2(p+d+1).
\]
The final auxiliary budget for label $i$ is $p+d\leq n_i-2$, so each
source has degree at most $2n_i-2$.  If
$\deg\omega_{i-1}\leq2n_{i-1}-1$, then
\[
 \deg\Phi^*\omega_{i-1}
 \leq2(2n_{i-1}-1)+1
 =4n_{i-1}-1\leq2n_i-1.
\]
Also $\deg dH_i=2n_i-1$.  These estimates prove every required degree
bound for~\eqref{eq:mechanical-global-form} and the final field.

At stage $i$, put $s=n_{i-1}$, $N=n_i=2s+\eta_i$. For an auxiliary
birth with $k$ remaining folds let $A=2^k$,
$d=\lfloor s/A\rfloor$, and $e=s-Ad=e_{i,k}$.  Its actual birth
half-degree is $d$.  The allowed modes in
\eqref{eq:mechanical-birth-modes} number
\[
 M_k=\sum_{q=0}^{d-1}\{N-A(2q+1)\}
     =d(N-Ad).
\]
All summands are positive: at the largest $q$, their value is
$2e+\eta_i+A\geq1$.  In particular $M_k\geq1$.  The $A$ real
copies contribute $A(M_k-1)$ simple zeros.  The exact identity
\[
 A M_k=(s-e)(s+\eta_i+e)
       =s(s+\eta_i)-e^2-\eta_i e
\]
shows that four source copies give
$4\sum_{k=0}^{i-1}A(M_k-1)=K_i$, as in
\eqref{eq:mechanical-birth}.  This also proves $K_i\geq0$ and
integrality. The selected-zero counts obey
$L_0=0$, $L_i=4L_{i-1}+K_i$, yielding
\eqref{eq:mechanical-total}.

Write $\omega_r=P\,dx+Q\,dy$.  Lemma~\ref{lem:melnikov}, applied to
the finite disjoint configuration, supplies one sufficiently small
nonzero $\mu$ for which
\[
 X_\mu=( (H_r)_y-\mu Q,\,-(H_r)_x+\mu P )
\]
has a hyperbolic limit cycle near each selected simple period zero.
This is one real polynomial vector field, of degree at most $2n-1$.
\end{proof}


\section{The finite estimate and all degrees}\label{sec:counting}

The construction gives $L_i=4L_{i-1}+K_i$, $L_0=0$, with
$K_i$ as in \eqref{eq:mechanical-birth}; hence \eqref{eq:mechanical-total}. We now control the
losses uniformly over all binary prefixes.
Put $s=n_{i-1}$ and $A_*=2^{i-1}$. Since $A_*\le s<2A_*$,
\begin{align}
 \sum_{k=0}^{i-1}e_{i,k}^2
 &\le (s-A_*)^2+\sum_{k=0}^{i-2}4^k
 \le(s-A_*)^2+\frac{A_*^2}{3}
 \le\frac{s^2}{3},\label{eq:residue-square}\\
 \sum_{k=0}^{i-1}e_{i,k}
 &\le 2^i-1-i.\label{eq:residue-linear}
\end{align}
The last inequality in \eqref{eq:residue-square} is equivalent to
$(s-A_*)(s-2A_*)\le0$. Define
\[
 B_i=i\eta_i+4\eta_i\sum_{k=0}^{i-1}e_{i,k}+4(2^i-1).
\]
Because $4s^2\le n_i^2$, equation \eqref{eq:mechanical-birth} implies
\begin{equation}\label{eq:countlower}
 L(n)\ge\sum_{i=1}^r(i-1/3)4^{r-i}n_i^2
             -\sum_{i=1}^r4^{r-i}B_i.
\end{equation}
The prefix identity $0\le n-2^{r-i}n_i<2^{r-i}$ gives
\begin{align*}
 \sum_{i=1}^r(i-1/3)4^{r-i}n_i^2
 &\ge n^2\left(\frac{r^2}{2}+\frac r6\right)
       -2n2^r\sum_{i=1}^{\infty}(i-1/3)2^{-i}\\
 &=n^2\left(\frac{r^2}{2}+\frac r6\right)-\frac{10}{3}n2^r.
\end{align*}
By \eqref{eq:residue-linear} and $0\le\eta_i\le1$,
\begin{align*}
 \sum_{i=1}^r4^{r-i}B_i
 &\le4^r\sum_{i=1}^{\infty}4^{-i}
       \{i+4(2^i-1-i)+4(2^i-1)\}\\
 &=4\cdot4^r.
\end{align*}
Substitution in \eqref{eq:countlower}, followed by $2^r\le n$, proves
\eqref{eq:finitebound}. If $n=2^r$, all remainders and all $\eta_i$
vanish. Summing $K_i=i4^i-4(2^i-1)$ gives \eqref{eq:dyadic}.

For any integer $D\ge3$ take $n=\lfloor(D+1)/2\rfloor$.
Then $2n-1\le D$, $n=D/2+O(1)$, and
$r=(\ln D)/(\ln2)+O(1)$. Equation \eqref{eq:finitebound} yields
\eqref{eq:asymptotic}. In particular the claimed coefficient is an
all-degree lower bound, rather than only a subsequence coefficient.
This completes the proof of Theorem~\ref{thm:main}.


\section{Comparison with the preceding constructions}\label{sec:replication}

\subsection{Two-source renewal and mechanical histories}\label{sec:scope}

Both the renewal theorem of \cite[Proposition~4.1]{Lu2026} and the
present construction use fourfold copying with two axis sources.
The earlier theorem starts from an
arbitrary polynomial field carrying a selected finite hyperbolic
configuration. It transplants that configuration into the interior of a
square and rebuilds the traces on two adjacent edges at a cost of two
degrees. A signed tangency condition makes all simple source roots
available for center creation. The classical fourfold transformation and
two successive common perturbations then give the update stated in the
introduction. The normalized number of new cycles per stage is of
constant order, producing the single logarithm.

Proposition~\ref{prop:mechanical-certificate} instead constructs the
separable Hamiltonians and perturbing forms together. Its source at
stage $i$ retains $i$ auxiliary birth levels, so the normalized new-zero
count is of order $i$. Summation over the planar stages supplies the
second logarithm. The future-degree allocation pays for this larger
source. Exact cancellation uses a reflection fixing each new oval and
reversing its orientation; exchanging two different ovals would only
make their periods opposite. Small even additions preserve this
cancellation while filling the intervening degrees.

Thus the two results address different input classes as well as
different quantitative bounds. The earlier renewal preserves a given
arbitrary hyperbolic configuration. The present theorem improves the
asymptotic count along an explicitly constructed near-Hamiltonian
sequence; it does not assert the same improved update for arbitrary
input fields. The proof constructs the auxiliary history and global
form directly, without the transplantation or signed edge traces used
in the earlier theorem.

\subsection{The Qin--Sun construction and all-degree comparison}

Qin--Sun's two logarithms arise from two 3-adic filtrations, one in each coordinate, in an
anisotropic Chebyshev center grid, with a simultaneous choice of
perturbation coefficients. Here they arise from the auxiliary mechanical
history and the stages of planar renewal. The square-logarithmic order
is therefore a prior result of Qin--Sun; the contribution here is the
mechanical construction and the improved coefficient. To compare the
coefficients over the same range of degrees, we first transfer their
subsequence bound to all sufficiently large degrees.

The required replication operation is due to
Eshkobilov--Kadyrov--Mamayusupov
\cite[Theorem~1, Lemma~2(4) and Section~6.3]{EKM2026}, whose
time-change argument also preserves hyperbolicity. We record the form
used here.

\begin{lemma}\label{lem:replication}
If a real polynomial field of degree at most $d$ has $N$ selected
hyperbolic limit cycles, then for every integer $m\ge1$ there is a
real polynomial field of degree at most $m(d+1)-1$ with at least
$m^2N$ hyperbolic limit cycles.
\end{lemma}
\begin{proof}
An affine change of coordinates places the selected compact cycles
inside $(-1,1)^2$. Write the resulting field as $(P,Q)$ and let $T_m$
be the first-kind Chebyshev polynomial. The polynomial field
\[
 \bigl(T_m'(y)P(T_m(x),T_m(y)),\,
       T_m'(x)Q(T_m(x),T_m(y))\bigr)
\]
has the stated degree. Each of the $m^2$ open monotonicity rectangles
maps diffeomorphically onto $(-1,1)^2$. Pushing the field forward gives
the seed field multiplied by the nonvanishing scalar
$T_m'(x)T_m'(y)$. Thus each selected cycle has one copy in every
rectangle. A positive time change preserves its return map, and a
negative time change replaces that map by its inverse. The multiplier
is consequently either the original multiplier or its reciprocal;
hyperbolicity persists in both cases.
\end{proof}

Combining the lemma with \eqref{eq:QS} gives the explicit bound
\begin{equation}\label{eq:envelope}
 \Hhyp(D)\ge B_{\mathrm{QS}}(D):=
 \max_{\substack{r\ge2\\4\cdot3^r-5\le D}}
 \left\lfloor\frac{D+1}{4\cdot3^r-4}\right\rfloor^2
 \bigl[4r^2 3^{2r-2}-(3^r-1)^2\bigr]
 \qquad(D\ge31).
\end{equation}
The case $m=1$ in Lemma~\ref{lem:replication} recovers each unreplicated
seed; the factor in \eqref{eq:envelope} chooses the largest admissible $m$
for each seed degree.

\begin{proposition}\label{prop:QSall}
For all sufficiently large integer $D$,
\begin{equation}\label{eq:QSalltransfer}
 \Hhyp(D)\ge\frac{D^2}{36(\ln3)^2}
 \bigl[(\ln D)^2-2\ln D\ln\ln D\bigr]-O(D^2\ln D).
\end{equation}
\end{proposition}
\begin{proof}
Put $S=D+1$ and choose $r$ maximal with
$4\cdot3^r\le S/\ln S$. Then $q=3^r=\Theta(S/\ln S)$ and
\[
 r=\frac{\ln D-\ln\ln D+O(1)}{\ln3}.
\]
Set $m=\lfloor S/(4q-4)\rfloor$. This meets the degree bound and
$mq=D/4+O(D/\ln D)$. Inserting these estimates in
$m^2[(4/9)q^2r^2-(q-1)^2]$ gives \eqref{eq:QSalltransfer}; the
remaining $(\ln\ln D)^2$ term is $O(\ln D)$.
\end{proof}

For clarity, the relevant coefficients are displayed in
Table~\ref{tab:comparison}. Qin--Sun's direct all-degree formula
\cite[Corollary~1.2 and Sections~6.2, 6.4]{QinSun2026} gives, for every $D\ge31$,
\[
 \Hhyp(D)\ge\frac{(D+5)^2}{324}
 \left(\frac{\ln((D+5)/12)}{\ln3}\right)^2
 -\frac{(D+1)^2}{16}.
\]
Its smaller coefficient comes from filling gaps between consecutive
ternary seed degrees by monotonicity. Proposition~\ref{prop:QSall}
is the consequence of their seeds and replication proved here, and is
not a restatement of that corollary.

\begin{table}[htbp]
\centering\small
\begin{tabular}{@{}p{.34\linewidth}p{.22\linewidth}p{.19\linewidth}p{.17\linewidth}@{}}
\toprule
Result & Degree range & Coefficient $c$ & Error factor\\
\midrule
Qin--Sun, Theorem 1.1 & $4\cdot3^r-5$, $r\ge2$ & $1/[36(\ln3)^2]$ & $\ln D$\\[3pt]
Qin--Sun, Corollary 1.2 & Every $D\ge31$ & $1/[324(\ln3)^2]$ & $\ln D$\\[3pt]
Qin--Sun + replication, Proposition~\ref{prop:QSall} & Every sufficiently large $D$ & $1/[36(\ln3)^2]$ & $\ln D\ln\ln D$\\[3pt]
Mechanical renewal, Theorem~\ref{thm:main} & Every sufficiently large $D$ & $1/[8(\ln2)^2]$ & $\ln D$\\
\bottomrule
\end{tabular}
\caption{Bounds of the form $\Hhyp(D)\ge cD^2(\ln D)^2$
with remainder $O(D^2)$ times the displayed error factor. All degrees are total degrees of one real
planar polynomial vector field; all counted cycles are hyperbolic.}
\label{tab:comparison}
\end{table}

\section{An iterated renewal conjecture}\label{sec:conjecture}

The preceding constructions suggest the following sufficient scheme for
higher logarithmic powers. Its unresolved step is the simultaneous
dynamical realization of new and inherited cycles at controlled degree.

\begin{conjecture}[Iterated renewal]\label{conj:renewal}
For each fixed integer $k\ge3$, there are integers $b,A\ge2$ and a
constant $c>0$, all depending only on $k$, with the following property.
For every integer $r\ge1$ there is a finite sequence of real planar
polynomial fields $X_0,\ldots,X_r$. For $0\le j\le r$,
$\deg X_j\le Ab^j-1$ and $X_j$ carries a selected finite family of $N_j$
hyperbolic limit cycles. At each step $1\le j\le r$, this family contains
the hyperbolic continuations of all $b^2$ pullback copies of each selected
cycle of $X_{j-1}$ and at least $c b^{2j}j^{k-1}$ additional cycles.
The pullback is separable of degree $b$, with $b$ real diffeomorphic
inverse branches in each coordinate over the preceding configuration.
All continued and additional cycles are pairwise distinct in $X_j$.
\end{conjecture}

The fields and auxiliary parameters may depend on the prescribed finite
horizon $r$, as in Proposition~\ref{prop:mechanical-certificate}. The input fields
are to be constructed as part of the sequence; the conjecture does not
require an update accepting every polynomial input field. A hyperbolic
continuation means that each selected pullback cycle persists hyperbolically
along a deformation from the pullback field to $X_j$ through polynomial
fields of degree at most $Ab^j-1$; its multiplier may vary.

The claimed construction would give
\[
 N_j\ge b^2N_{j-1}+c b^{2j}j^{k-1},\qquad
 \frac{N_r}{b^{2r}}\ge c\sum_{j=1}^rj^{k-1}\ge\frac{c}{k}r^k.
\]
For large $D$, taking $r=\lfloor\log_b((D+1)/A)\rfloor$ gives $Ab^r-1\le D$,
$b^r>(D+1)/(Ab)$ and $r\sim\ln D/\ln b$. Thus the conjecture would imply
\[
 \Hhyp(D)\ge c_k D^2(\ln D)^k
\]
for some $c_k>0$ and all sufficiently large $D$. Degreewise existence of
a field with a lower logarithmic power does not itself provide the
renewal asserted in Conjecture~\ref{conj:renewal}: the new cycles must be
created together with the retained copies at every stage without an
additional degree cost that destroys the estimate. The constructions
for one and two logarithms motivate this question; no such iterated
realization for $k\ge3$ is proved here.


\section*{Reproducibility and use of AI}

The accompanying program \texttt{repro/verify.py} checks the finite
counts, degree inequalities and replication envelope by integer
arithmetic. These checks supplement the analytic proofs above.

OpenAI's ChatGPT and Codex were used to explore proof strategies,
check calculations, develop verification scripts, and revise the text.
The author is responsible for all mathematical claims and the final
manuscript.

\bibliographystyle{plain}
\bibliography{references}
\end{document}

